\documentclass[11pt]{letter}

\usepackage[margin=1in]{geometry}
\usepackage{amsmath}

\signature{Stephen Nelson\\Field Dynamics Research}

\begin{document}

\begin{letter}{%
The Editor\\
\textit{Classical and Quantum Gravity}%
}

\opening{Dear Editor,}

I am submitting the enclosed letter, ``The Coordinate-Space Geometry of the Tetrad Frame,'' for consideration in \textit{Classical and Quantum Gravity}.

The tetrad has been the standard variable for coupling gravity to fermionic matter since the work of Cartan, Einstein, and Weyl. It is central to Poincar\'{e} gauge gravity, Ashtekar variables, loop quantum gravity, bimetric gravity, and supergravity. In over a century of use, the spatial geometry generated by its local frame has not, to our knowledge, been explicitly derived in coordinate space.

This submission identifies a previously unreported coordinate-space geometry associated with the local tetrad frame. Starting from a single orthonormal tetrad, the paper shows that the induced half-space structure yields exactly two bounded tetrahedral regions whose union forms a stellated octahedron with a regular octahedral core. The opposite-oriented tetrahedron is not introduced independently; it emerges as the unique second bounded region of the same construction. Because the result is kinematic and tetrad-level, it arises prior to any choice of gravitational action or field equation and is therefore relevant across multiple tetrad-based formalisms.

The geometry is established by two complementary derivations: a convex boundedness classification in the main letter, and a supplementary group-theoretic derivation from the quadratic tetrad-to-metric relation. Supporting materials also include a supplementary note deriving the tetrahedral frame from the Lorentz signature's isotropy, an interactive rendering, and a validation script confirming the principal geometric claims.

I believe referees with expertise in tetrad formalism, vierbein geometry, or Poincar\'{e} gauge gravity would be well placed to evaluate this work.

\closing{Yours sincerely,}

\end{letter}

\end{document}
